% ISSRE 2026 Fast Abstract
% Track: Fast Abstract & Project Highlights
% Deadline: June 15, 2026
% Format: IEEE Computer Society, 2 pages max
% Compile: pdflatex issre2026-fast-abstract.tex && bibtex issre2026-fast-abstract && pdflatex issre2026-fast-abstract.tex && pdflatex issre2026-fast-abstract.tex

\documentclass[conference,10pt]{IEEEtran}

\usepackage[utf8]{inputenc}
\usepackage[T1]{fontenc}
\usepackage{amsmath,amssymb}
\usepackage{booktabs}
\usepackage{hyperref}
\usepackage{url}
\usepackage{enumitem}

\title{FaultRay: In-Memory Cascade Simulation\\with Formal LTS Semantics and\\Multi-Layer Availability Limits}

\author{
  \IEEEauthorblockN{Yutaro Maeda}
  \IEEEauthorblockA{Independent Researcher\\
    Chigasaki, Japan\\
    contact@faultray.com}
}

\begin{document}
\maketitle

%% ---- ABSTRACT ----
\begin{abstract}
Production chaos engineering tools inject real faults into running
systems, creating risk for regulated operators. Classical reliability
analysis (FTA, RBD) assumes component independence, which does not hold
for cloud infrastructure where a single incident can simultaneously
degrade multiple layers. We present \textsc{FaultRay}, a research
prototype contributing: (1)~a cascade propagation engine formalized as
a Labeled Transition System (LTS) over a dependency graph, with proven
termination and $O(|V|{+}|E|)$ complexity; and (2)~an $N$-layer
availability limit model composing categorical ceilings via
$\min$ rather than the classical series-product. An illustrative
backtest on 54~real-world cloud incidents reproduces known cascade
paths. \textsc{FaultRay} is open source (Apache~2.0, PyPI).
\end{abstract}

%% ---- 1. PROBLEM ----
\section{Problem and Motivation}

Modern cloud infrastructure cascading failures---the 2021 Meta BGP
outage, the 2017 AWS S3 incident, the 2022 Cloudflare BGP
leak---demonstrate how a single fault propagates across dozens of
dependent services within minutes. Two gaps remain in current
approaches:

\textbf{Gap 1: Production risk.} Chaos engineering tools (Gremlin,
Steadybit, AWS FIS, LitmusChaos) inject \emph{real} faults. For
regulated financial institutions under DORA (EU~2022/2554), touching
production to discover weaknesses is often unacceptable.

\textbf{Gap 2: Independence assumption.} Classical Reliability Block
Diagrams compose availabilities via series-product
$A = \prod_i A_i$~\cite{trivedi2001}, assuming independent failures.
Cloud infrastructure violates this: a region outage simultaneously
degrades compute, network, and storage layers.

Concurrent work by Krasnovsky~\cite{krasnovsky2026} (ICSE-NIER~2026)
proposes in-memory Monte-Carlo simulation on service-dependency
graphs. \textsc{FaultRay} addresses the same positioning with a
different technical approach: formal LTS semantics with proven
properties and multi-layer availability decomposition.

%% ---- 2. APPROACH ----
\section{Approach}

\subsection{Cascade Engine (LTS Formalization)}

\textsc{FaultRay} models cascade propagation as a Labeled Transition
System $\mathcal{M} = (S, s_0, \mathit{Act}, {\longrightarrow}, F)$
over an infrastructure dependency graph $G = (C, E, \mathit{dep},
w, \tau, \mathit{cb})$. Each state is a 4-tuple $(H, L, T, V)$:
health map, latency map, elapsed time, and visited set.

Eight transition rules govern fault injection, propagation through
typed dependencies (\texttt{requires}/\texttt{optional}/\texttt{async}),
circuit breaker trips, timeout cascades, and termination. Proven
properties include:
\begin{itemize}[nosep,leftmargin=*]
  \item \textbf{Termination}: the visited set $V$ grows monotonically
    and is bounded by $|C|$.
  \item \textbf{$O(|V|{+}|E|)$ complexity}: each component is visited
    at most once; each edge is examined once.
  \item \textbf{Monotonicity}: health can only worsen during a run.
  \item \textbf{Bounded blast radius}: affected components are bounded
    by the reverse-reachable set from the fault origin.
\end{itemize}

\subsection{$N$-Layer Availability Limit Model}

Rather than a single $A = \text{MTBF}/(\text{MTBF}+\text{MTTR})$,
\textsc{FaultRay} decomposes availability into five categorical layers:
\begin{equation}
  A_{\text{sys}} = \min(A_{\text{L1}}, A_{\text{L2}}, A_{\text{L3}},
  A_{\text{L4}}, A_{\text{L5}})
\end{equation}
where L1=software (deploy, drift), L2=hardware (MTBF/MTTR, redundancy),
L3=theoretical (packet loss, GC), L4=operational (incident response),
L5=external SLA ($\prod a_i$). The $\min$ operator replaces the
series-product because cloud layers exhibit \emph{correlated} failure
modes, making the weakest-link composition more appropriate than
independent-product composition.

%% ---- 3. PRELIMINARY RESULTS ----
\section{Preliminary Results}

\textbf{Backtest.} We performed a post-hoc backtest against
54~real-world cloud incidents (2017--2024: AWS, GCP, Azure, Meta,
Cloudflare, GitHub, and 16 others). For each, a topology graph was
constructed from the public post-mortem and the documented root
cause was injected. The cascade engine achieved $F_1 = 0.87$ (precision $= 1.00$,
recall $= 0.81$, severity accuracy $= 0.58$). We stress this is
\emph{illustrative}, not predictive: topologies were constructed
with knowledge of the outcome.

\textbf{Scale.} The engine processes a 50-component topology
in ${<}10$\,ms on commodity hardware (Apple M1). The implementation
comprises 512~source files and 33,841~tests.

\textbf{Positioning.} Table~\ref{tab:pos} summarizes key differences
from representative tools and concurrent work.

\begin{table}[t]
\centering
\caption{Positioning of \textsc{FaultRay}.}
\label{tab:pos}
\footnotesize
\begin{tabular}{@{}lccc@{}}
\toprule
& \textbf{FaultRay} & \textbf{Krasnovsky} & \textbf{Chaos tools} \\
\midrule
In-memory       & \checkmark & \checkmark & \texttimes \\
Formal proofs   & \checkmark & \texttimes & \texttimes \\
$N$-layer avail.& \checkmark & \texttimes & \texttimes \\
Monte-Carlo     & \texttimes & \checkmark & \texttimes \\
Real injection  & \texttimes & \texttimes & \checkmark \\
Open source     & \checkmark & \checkmark & partial \\
\bottomrule
\end{tabular}
\end{table}

%% ---- 4. STATUS AND FUTURE ----
\section{Status and Future Work}

\textsc{FaultRay} is released as open source under Apache~2.0
on GitHub\footnote{\url{https://github.com/mattyopon/faultray}}
and PyPI. A US provisional patent has been filed
(No.~64/010,200, 2026-03-19). A Zenodo preprint is available
(DOI:~10.5281/zenodo.19139911).

Future work includes: (1)~prospective validation on live
infrastructure with blinded topologies; (2)~integration with
DeathStarBench~\cite{gan2019} for quantitative comparison with
Krasnovsky's Monte-Carlo approach; (3)~extending the $N$-layer
model to capture temporal correlation between layers.

\textbf{Disclaimer.} \textsc{FaultRay} is a research prototype,
not a certified compliance tool.

%% ---- REFERENCES ----
\bibliographystyle{IEEEtran}
\begin{thebibliography}{9}

\bibitem{trivedi2001}
K.~S.~Trivedi,
\textit{Probability and Statistics with Reliability, Queuing, and
Computer Science Applications}, 2nd~ed.
Wiley, 2001.

\bibitem{krasnovsky2026}
A.~A.~Krasnovsky,
``Model Discovery and Graph Simulation: A Lightweight Gateway to
Chaos Engineering,''
in \textit{Proc.\ ICSE-NIER}, 2026.
arXiv:2506.11176.

\bibitem{gan2019}
Y.~Gan \textit{et al.},
``An Open-Source Benchmark Suite for Microservices and Their
Hardware-Software Implications for Cloud \& Edge Systems,''
in \textit{Proc.\ ASPLOS}, 2019.

\end{thebibliography}

\end{document}
